\RequirePackage[T1]{fontenc}
\documentclass[letterpaper,10pt,conference]{ieeeconf}
\overrideIEEEmargins
\usepackage{amsmath,amssymb}
\usepackage{url}
\title{Temporal State-Space Adapters in Frozen Image Encoders}
\author{Working draft: authors and affiliations to be supplied}
\begin{document}
\maketitle
\begin{abstract}
We investigate whether interleaving recurrent temporal updates with a frozen
image encoder improves the accuracy--cost trade-off of streaming visual
recognition. The implementation combines SigLIP 2 spatial processing with
per-coordinate Mamba-3 adapters and exposes only current-frame tokens to a
common task readout. The study is designed to compare image-only, temporal
attention, late recurrent processing, and published GatedMambaMix cores.
This working draft reports an intermediate learning curve and a completed
history intervention diagnostic. Matched final-checkpoint comparisons,
multiple-seed results, and long-range recognition evidence remain pending;
the present results do not establish architectural superiority.
\end{abstract}

\section{Research Question and Scope}
Robotic perception benefits from distinguishing observations with similar
current appearances but different histories. Our target is a visual encoder
that uses those histories without expanding the downstream token sequence.
We do not modify a robot policy or infer manipulation success from video
classification. The primary question is whether recurrent updates within a
strong spatial encoder offer a useful recognition--cost trade-off compared
with temporal attention under matched data and training conditions.

The current implementation is a research baseline. Using Mamba-3, inserting
time processing between spatial layers, and returning current-frame tokens
are not individually novel contributions. A publishable performance claim
requires completed comparisons, a mechanism explaining the improvement,
and a robotics-relevant recognition evaluation. Those requirements are
not replaced by implementation correctness or the diagnostic below.

\section{Related Work}
SigLIP 2 provides pretrained visual representations with strong semantic
and dense-task capabilities~\cite{siglip2}. We reuse its spatial weights.
Mamba-3 develops state-space recurrences for sequence modeling~\cite{mamba3};
its language-model results do not imply a visual-recognition advantage.

MEM uses temporal attention at regular ViT depths and forwards current
tokens to a VLA, while handling longer semantic history with a separate
language memory~\cite{mem}. TRecViT alternates recurrent temporal mixing,
spatial attention, and channel processing~\cite{trecvit}. SDQES inserts
streaming temporal adapters into image backbones~\cite{sdqes}. These works
preclude a broad novelty claim for recurrent image-to-video adaptation.

Frozen Image + GMMix evaluates temporal cores on frozen image features,
including multi-depth features~\cite{gmmix}. We include both final-layer
and multi-depth variants of its public core. Replacing its backbone with
SigLIP 2 is a controlled adaptation, not a reproduction of the paper's
reported DINO results. Our temporal-attention control similarly does not
reproduce MEM's shared-parameter attention or its robot training.

\section{Encoder and Learning}
Let $I_t$ be the current image and $X_t^\ell\in\mathbb{R}^{N\times D}$
the features after spatial block $\ell$. For SigLIP 2-B/16 at $224^2$
resolution, $N=196$ and $D=768$. We insert adapters after blocks 4, 8,
and 12. At patch coordinate $n$, each adapter computes
\begin{align}
 u_{t,n}^\ell &= W_\downarrow^\ell
       \operatorname{LN}^\ell(X_{t,n}^\ell)+b_\downarrow^\ell,\\
 (v_{t,n}^\ell,s_{t,n}^\ell) &=
       \mathcal{M}_\theta^\ell(u_{t,n}^\ell,s_{t-1,n}^\ell),\\
 \widetilde X_{t,n}^\ell &= X_{t,n}^\ell+
       g^\ell\odot(W_\uparrow^\ell v_{t,n}^\ell+b_\uparrow^\ell).
\end{align}
Here $u$ has width 256. $\mathcal M$ is the upstream Mamba-3 SISO core
with state dimension 64, expansion 2, and head dimension 64. The channel
gate $g$ starts at zero. Each layer has independent learned parameters.
At initialization the residual update is zero; this is not a guarantee
that the pretrained spatial behavior remains unchanged after training.

Temporal recurrence is independent at each image coordinate; subsequent
SigLIP spatial attention mixes coordinates. There is no explicit object
tracking or motion-aligned memory addressing. The native recurrent cache
includes accumulated angles, the SSM matrix, and previous key/value
states. All four are carried between observations and reset between
episodes. Cameras are separate batch elements in the current interface.

After final layer normalization, the encoder emits
$Z_t\in\mathbb{R}^{196\times768}$. A common attentive readout predicts
the action label from $Z_T$ only. All past observations pass through the
spatial and temporal computation; they are not separately passed to the
readout. The image-only baseline processes the same final image.

We freeze SigLIP weights and train the adapters and readout with label
cross-entropy, without distillation. Gradients still pass through later
frozen spatial blocks to earlier adapters. Freezing parameters therefore
does not remove all backbone backward computation. Clip training starts
from empty memory. Inference supports a persistent native cache.

\subsection{Capacity and Cost}
At fixed resolution, architecture, and batch size, recurrent state capacity
does not depend on history length. This property concerns inference state,
not training activation memory or information retention. The measured
native cache for this architecture at batch one is 79,027,200 bytes.
Our three-layer width-256 BF16 attention control stores
\begin{equation}
 M_{\mathrm{KV}}(T)=3\cdot196\cdot2\cdot256\cdot2\cdot T
                  =602{,}112T\ \mathrm{bytes}.
\end{equation}
Consequently its cache is smaller for short sequences; recurrent cache
capacity becomes smaller at 132 observations. This arithmetic is not a
crossover claim for total GPU memory or latency. Spatial encoder weights,
compute, and the 196 output tokens are retained by both methods.

\section{Experimental Protocol}
\subsection{Matched Recognition Comparisons}
The fixed protocol uses 168,913 training and 24,777 validation SSv2 videos
with 174 classes. Eight sparse observations are sampled over each clip,
including the same last observation; they are not eight consecutive frames.
Training uses shared spatial crops, no horizontal flips, AdamW with
learning rate $10^{-4}$ and weight decay $10^{-4}$, 1,000 warmup updates,
and 20,000 total updates. Microbatch four with accumulation eight gives
effective batch 32. Spatial computation is FP32 with TF32 disabled;
temporal computation uses BF16. Comparisons share the backbone revision,
data, readout family, and final-update selection rule.

The six-way seed-42 grid consists of image-only SigLIP 2, interleaved
temporal attention, interleaved Mamba-3, late Mamba-3, and two GMMix
feature configurations. Interleaved and late Mamba use identical temporal
parameter counts. Other architectures are not parameter matched, so
accuracy must be assessed with measured cost and capacity. Major contrasts
are scheduled to be repeated with seeds 43 and 44. Three seeds provide limited uncertainty
information. Validation has been inspected during development and is not
an untouched test set.

The primary metric is full-history streaming top-1. The planned analysis reports paired
prediction differences, conditional source-video bootstrap intervals,
and training-seed variation separately. GMMix uses the documented
PyTorch convolution fallback on this machine; kernel differences limit
hardware-general speed conclusions.

\begin{table}[t]
\centering\small
\caption{Planned final-checkpoint comparison. Dashes indicate results
not yet supplied by completed matched runs, not zero accuracy.}
\begin{tabular}{lc}
\hline
Final comparison, 20,000 updates & Streaming top-1 \\
\hline
Image only & -- \\
Interleaved temporal attention & -- \\
Interleaved Mamba-3 & -- \\
Late Mamba-3 & -- \\
Final-layer GMMix & -- \\
Multi-depth GMMix & -- \\
\hline
\end{tabular}

\end{table}

\subsection{Streaming Costs and Long-Range Recognition}
Cost measurement uses the same fixed real RGB sequences of 52, 245, and
1,337 observations, batch one, one warmup traversal, and three measured
traversals. GPU latency, preprocessing/transfer-inclusive latency, state
capacity, and peak allocated memory are distinct metrics. Timing runs
exclude concurrent GPU jobs. File decoding is outside the timed path.
These long sequences measure computational scaling, not long-term memory
accuracy. Recognition claims additionally require evidence-dependent
queries with a known evidence-to-query distance and valid target labels.
That evaluation is not yet completed.

\section{Completed Intermediate Evidence}
\begin{table}[t]
\centering\small
\caption{Intermediate clip validation, one seed and all 24,777 videos.
These checkpoints are not the final comparison checkpoints.}
\begin{tabular}{rrr}
\hline
Updates & Correct / 24,777 & Top-1 (\%) \\
\hline
1,000 & 3,990 & 16.10 \\
2,000 & 7,203 & 29.07 \\
3,000 & 8,568 & 34.58 \\
4,000 & 9,710 & 39.19 \\
5,000 & 10,879 & 43.91 \\
6,000 & 11,448 & 46.20 \\
7,000 & 11,875 & 47.93 \\
8,000 & 12,453 & 50.26 \\
9,000 & 12,551 & 50.66 \\
\hline
\end{tabular}

\end{table}
The learning curve is regenerated from a saved audit of official labels,
data splits, source hashes, and validation prediction counts. It is not
an independent rerun of every GPU forward. A corrected RGB conversion is
used throughout; results before that correction are excluded.

\subsection{Fixed-Current-Image History Interventions}
We use one archived 4,000-update checkpoint and two validation videos per
class sampled with seed 17. The resulting 348-video set was inspected in
earlier diagnostics. We keep weights, classifier, and the last image fixed,
and alter only history. Current-only inference uses a fresh state; the
repeated-current condition uses eight identical current observations;
history reversal reverses only the first seven observations.

\begin{table}[t]
\centering\small
\caption{History intervention at 4,000 updates on the same 348 videos.
Scores retain original clip labels even under altered histories.}
\begin{tabular}{lrr}
\hline
History intervention & Correct / 348 & Top-1 (\%) \\
\hline
Original history, clip & 105 & 30.17 \\
Original history, stream & 105 & 30.17 \\
Current image only & 26 & 7.47 \\
Current image repeated & 27 & 7.76 \\
Past observations reversed & 93 & 26.72 \\
\hline
\end{tabular}

\end{table}
The original-history condition answers more labels correctly than either
current-only condition. Repeating the current observation does not recover
the original-history performance. This supports history use by the fitted
model. It does not compare independently trained architectures. Reversing
history can change the action's meaning, so the original-label difference
is not an accuracy measure for relabelled reversed videos.

Original clip logits exactly reproduce the prior saved diagnostic.
Clip and streaming logits differ by at most 0.0625, with no predicted-class
disagreements on these 348 examples. All 1,740 condition/example predictions
were recounted against the original labels. No conclusion about longer
sequences follows from this short-clip agreement.

\section{Outstanding Evidence and Decision Criteria}
Final six-way accuracy/cost results and the three-seed comparisons must be
completed before making a performance claim. The main contrast is
interleaved Mamba versus temporal attention; interleaved versus late Mamba
isolates placement. A positive result must not be explained solely by
additional parameters, different kernels, data, or a changed readout.

Robotics relevance requires a history-dependent perception application
with controlled current observations or evidence timing. Fine point
correspondence is not the main success criterion. Previous geometric
pilots are auxiliary diagnostics and are not evidence of robot success.
No action policy or physical robot evaluation has been performed in this
study. If recognition and cost comparisons do not show a reproducible
advantage, the architectural hypothesis must be revised or rejected.

\begin{thebibliography}{9}
\bibitem{siglip2} M. Tschannen et al., ``SigLIP 2: Multilingual
Vision-Language Encoders with Improved Semantic Understanding,
Localization, and Dense Features,'' arXiv:2502.14786, 2025.
\bibitem{mamba3} A. Lahoti et al., ``Mamba-3: Improved Sequence Modeling
using State Space Principles,'' arXiv:2603.15569, 2026.
\bibitem{mem} M. Torne et al., ``MEM: Multi-Scale Embodied Memory for
Vision Language Action Models,'' arXiv:2603.03596, 2026.
\bibitem{trecvit} V. P\u{a}tr\u{a}ucean et al., ``TRecViT: A Recurrent
Video Transformer,'' arXiv:2412.14294, 2024, revised 2026.
\bibitem{sdqes} C. Eyzaguirre et al., ``Streaming Detection of Queried
Event Start,'' arXiv:2412.03567, 2024.
\bibitem{gmmix} S. Orlova, N. Cavagnero, and G. Dubbelman,
``Towards Data-Efficient Video Pre-training with Frozen Image Foundation
Models,'' arXiv:2605.19137, 2026.
\end{thebibliography}
\end{document}
